% Zone „Das kennst du schon“ – aus bank/lineare-gleichungen/zone.jsonl (zusammenbau v0.1)
\einheitenkopf[zone][Kennst du schon]{Das kennst du schon}
\setcounter{aufgabe}{0}
%% TODO Grundvorstellungs-Aufgabe als erste Hauptnummer der Zone (2.8) fehlt in der Bank

%% TODO Titel: Ich-kann-Satz fehlt in der Bank (Platzhalter: Kettenname) – Nr. 1
%% TODO Anweisung: ein Satz über der ersten Teilaufgabe fehlt in der Bank – Nr. 1
\begin{aufgabe}{Punkt vor Strich (3 · 4 − 5)}
%% TODO Darstellung für schwach fehlt in der Bank (lineare-gleichungen-zone-f1-v1; Abschnitt „Für schwache Schüler“)
\swz{$2 + 4 \cdot 5$}{}
%% TODO Darstellung für schwach fehlt in der Bank (lineare-gleichungen-zone-f1-v3; Abschnitt „Für schwache Schüler“)
\swz{$4 \cdot (-3) + 20$}{}
\end{aufgabe}

%% TODO Titel: Ich-kann-Satz fehlt in der Bank (Platzhalter: Kettenname) – Nr. 2
%% TODO Anweisung: ein Satz über der ersten Teilaufgabe fehlt in der Bank – Nr. 2
\begin{aufgabe}{Negative Zahlen addieren, subtrahieren, dividieren (12 : (−3))}
%% TODO Darstellung für schwach fehlt in der Bank (lineare-gleichungen-zone-f2-v1; Abschnitt „Für schwache Schüler“)
\swz{$-3 + 8$}{}
%% TODO Darstellung für schwach fehlt in der Bank (lineare-gleichungen-zone-f2-v3; Abschnitt „Für schwache Schüler“)
\swz{$-2{,}5 - 1{,}5$}{}
\end{aufgabe}

%% TODO Titel: Ich-kann-Satz fehlt in der Bank (Platzhalter: Kettenname) – Nr. 3
%% TODO Anweisung: ein Satz über der ersten Teilaufgabe fehlt in der Bank – Nr. 3
\begin{aufgabe}{Terme zusammenfassen (gleichartige Glieder mit x, Zahl und x gemischt)}
%% TODO Darstellung für schwach fehlt in der Bank (lineare-gleichungen-zone-f3-v1; Abschnitt „Für schwache Schüler“)
\swz{$4x + 3x$}{}
%% TODO Darstellung für schwach fehlt in der Bank (lineare-gleichungen-zone-f3-v3; Abschnitt „Für schwache Schüler“)
\swz{$2x + 5 + 3x - 1$}{}
\end{aufgabe}

%% TODO Titel: Ich-kann-Satz fehlt in der Bank (Platzhalter: Kettenname) – Nr. 4
%% TODO Anweisung: ein Satz über der ersten Teilaufgabe fehlt in der Bank – Nr. 4
\begin{aufgabe}{Brüche: Hauptnenner, Bruch mal Zahl}
%% TODO Darstellung für schwach fehlt in der Bank (lineare-gleichungen-zone-f4-v1; Abschnitt „Für schwache Schüler“)
\swz{$\frac{1}{3} \cdot 12$}{}
%% TODO Darstellung für schwach fehlt in der Bank (lineare-gleichungen-zone-f4-v3; Abschnitt „Für schwache Schüler“)
\swz{$\frac{3}{4} \cdot 20$}{}
\end{aufgabe}

%% TODO Titel: Ich-kann-Satz fehlt in der Bank (Platzhalter: Kettenname) – Nr. 5
%% TODO Anweisung: ein Satz über der ersten Teilaufgabe fehlt in der Bank – Nr. 5
\begin{aufgabe}{Klammern auflösen}
%% TODO Darstellung für schwach fehlt in der Bank (lineare-gleichungen-zone-f5-v1; Abschnitt „Für schwache Schüler“)
\swz{$4 \cdot (x + 3)$}{}
%% TODO Darstellung für schwach fehlt in der Bank (lineare-gleichungen-zone-f5-v3; Abschnitt „Für schwache Schüler“)
\swz{$-3 \cdot (x - 4)$}{}
\end{aufgabe}

%% TODO Titel: Ich-kann-Satz fehlt in der Bank (Platzhalter: Kettenname) – Nr. 6
%% TODO Anweisung: ein Satz über der ersten Teilaufgabe fehlt in der Bank – Nr. 6
\begin{aufgabe}{Terme zusammenfassen (gleichartige Glieder mit x, Zahl und x gemischt) – Fehler finden}
\swz[3]{Tom fasst zusammen: $3x + 5 + 4x = 12x$. Was ist falsch?}{}
\end{aufgabe}

%% TODO Titel: Ich-kann-Satz fehlt in der Bank (Platzhalter: Kettenname) – Nr. 7
%% TODO Anweisung: ein Satz über der ersten Teilaufgabe fehlt in der Bank – Nr. 7
\begin{aufgabe}{Terme zusammenfassen (gleichartige Glieder mit x, Zahl und x gemischt) – selbst rechnen}
%% TODO Darstellung für schwach fehlt in der Bank (lineare-gleichungen-zone-f3-v8; Abschnitt „Für schwache Schüler“)
\swz{$6x + 4 + 3x$}{}
\end{aufgabe}
